% 共通プリアンブル（LuaLaTeX + 日本語）
% ビルド: lualatex 対象ファイル.tex（2回推奨）

\usepackage[hiragino-pron,deluxe]{luatexja-preset}
\usepackage{geometry}
\geometry{left=25mm,right=25mm,top=25mm,bottom=28mm}

\usepackage{graphicx}
\usepackage{booktabs}
\usepackage{longtable}
\usepackage{array}
\usepackage{hyperref}
\usepackage{xcolor}
\usepackage{listings}
\usepackage{fancyhdr}
\usepackage{enumitem}
\usepackage{tabularx}
\usepackage{fancyvrb}

\hypersetup{
  colorlinks=true,
  linkcolor=teal!60!black,
  urlcolor=teal!50!black,
  pdfborder={0 0 0},
}

% フォントは luatexja-preset（hiragino-pron）で設定済み
% Linux 等では common.tex 先頭を [noto-otf,deluxe] に変更してください
\setsansjfont{Hiragino Sans}
\setmonofont{Menlo}[Scale = 0.88]

\lstset{
  basicstyle = \ttfamily\small,
  breaklines = true,
  frame = single,
  rulecolor = \color{black!15},
  backgroundcolor = \color{black!3},
  columns = fullflexible,
  keepspaces = true,
  showstringspaces = false,
}

\newcommand{\doctypemark}{}
\newcommand{\setdoctype}[1]{\renewcommand{\doctypemark}{#1}}

\pagestyle{fancy}
\fancyhf{}
\fancyhead[L]{\small\leftmark}
\fancyhead[R]{\small\doctypemark}
\fancyfoot[C]{\thepage}
\renewcommand{\headrulewidth}{0.4pt}
\renewcommand{\footrulewidth}{0pt}

\setlist[itemize]{topsep=0.4em,itemsep=0.2em}
\setlist[enumerate]{topsep=0.4em,itemsep=0.2em}

\DefineVerbatimEnvironment{asciibox}{Verbatim}{%
  fontsize=\small,
  frame=single,
  rulecolor=\color{black!20},
  framesep=6pt,
  xleftmargin=1em,
  xrightmargin=1em,
}

\newcommand{\pkg}[1]{\texttt{#1}}
\newcommand{\code}[1]{\texttt{#1}}
